% SOURCE docs/release/2026-10-29/source/docs/writing-repair/2026-10-29/src/73ae9588bcea9fe8.tex CHAR 0
\chapter{Controlled-state identification, estimation and customer choice}
\label{r3:controlled-identification}
\section{What a hidden-state fit must establish}
The likelihood and policy-replacement recursions in
Chapter~\ref{r3:longitudinal} leave an identification question open.
Which controlled transition, reward and observation kernels can the
recorded law determine? The answer depends on who sees the state,
how actions are assigned, what is measured, and which actions the target
policy will use. A repayment trajectory can identify a useful predictive
state without identifying the effect of an untried restructuring.

For an uncontrolled stationary $K$-state HMM,
\citet[Assumptions A1--A3, Theorem 1]{hmmident2016} identifies a
full-rank ergodic transition matrix and distinct emission distributions
from a block of $2K+1$ observations, up to state-label permutation.
Distinct emissions need not be linearly independent in that theorem.
It is a result about the specified uncontrolled observation law.
Treating a privately selected action as an exogenous input changes its
factorization, so the theorem cannot simply be transferred to that fit.

Identification of the value of one supported policy can require less
information than identification of every state parameter. Conversely,
identifying all parameters of one chosen observational model says
nothing about excluded causal models with the same observed law.
The following finite construction separates these issues explicitly.

\section{A constructive controlled identification result}
\label{r3:known-emission}
Let $S_t\in\{1,\ldots,K\}$, with column vector of probabilities
$\mu_t$, and let $O_t\in\{1,\ldots,m\}$. For this subsection,
the predecision emission matrix $E$ is known from a justified validation
model, $E_{os}=P(O_t=o\mid S_t=s)$. Assume it has full column rank,
and let $L$ be a left inverse, $LE=I_K$. Assume emissions at adjacent
dates are independent conditional on their states, invariant across
actions, and measured under the same observation scheme.

At date $t$, the action is randomized independently of the state and
predecision emission, with positive probability for each action of
interest. Write $T_a(s,s')=P(S_{t+1}=s'\mid S_t=s,\doop(a))$,
with row-stochastic $T_a$, and assume the controlled structural model
makes the experiment identify this kernel. The observable joint matrix
conditional on randomized action satisfies
\begin{equation}
 J_a(o,o')=P(O_t=o,O_{t+1}=o'\mid A_t=a)
       =\{E\,\operatorname{diag}(\mu_t)T_aE'\}_{oo'}.
 \label{r3:emission-matrix}
\end{equation}
The predecision emission distribution gives $\mu_t=Lp_{O_t}$.
Left- and right-multiplying Equation~\eqref{r3:emission-matrix} gives
\begin{equation}
 T_a=\operatorname{diag}(\mu_t)^{-1}LJ_aL',
 \qquad \mu_t(s)>0\ \text{for every recovered row}.
 \label{r3:controlled-inversion}
\end{equation}
This is a constructive sufficient result under known emissions and
randomized actions, not an unknown-emission POMDP identification theorem.
If emission distributions depend on the action, the appropriate known
pre- and postdecision matrices and their inverses must replace $E$.

For a reward independent of $O_t$ conditional on $(S_t,A_t)$, let
$r_a(s)=\E[R_t\mid S_t=s,\doop(a)]$ and
$k_a(o)=\E[R_t\mathbf1(O_t=o)\mid A_t=a]$. Then
$k_a=E\operatorname{diag}(\mu_t)r_a$ and
$r_a=\operatorname{diag}(\mu_t)^{-1}Lk_a$.
The experiment, measurement model and state/action support together
identify these mean-reward and transition objects. None of the three
can be replaced by a good likelihood score.

\subsection{An observed reward can carry additional state information}
\label{r3:marked-reward-identification}
Mean rewards suffice for evaluating a policy that ignores reward as
information and uses the recovered signal process. They need not suffice
for the full-information belief recursion in
Equation~\eqref{r3:controlled-filter}. If the bank observes the reward
before its next decision, its conditional distribution can change the
posterior over states even when its mean is unchanged.

To identify that richer recursion, strengthen the emission restriction:
conditional on $(S_t,S_{t+1},A_t,R_t)$ the two signals have independent
laws $E_{\cdot S_t}$ and $E_{\cdot S_{t+1}}$. For every measurable
reward set $D$, the observed marked joint matrix then satisfies
\[
 J_{a,D}(o,o')
 =P(O_t=o,R_t\in D,O_{t+1}=o'\mid A_t=a)
 =E\operatorname{diag}(\mu_t)K_{a,D}E',
\]
where
$K_{a,D}(s,s')=T_a(s,s')P(R_t\in D\mid s,a,s')$.
The same multiplication yields
\[
 K_{a,D}=\operatorname{diag}(\mu_t)^{-1}LJ_{a,D}L'.
\]
Thus the marked observed law identifies the transition--reward measure,
and division by $T_a(s,s')>0$ identifies its reward conditional law
on supported state transitions. For finite reward categories this is
one matrix inversion per category. For general rewards it is an
identification statement about all measurable $D$, not a finite-sample
density-estimation guarantee. Together with known invariant emissions,
these joint kernels supply the information update as well as expected
reward. Merely recovering $r_a$ does not establish this stronger result.

For a summary-preservation counterexample, let a fair binary state persist,
with a binary signal correct with probability $4/5$. In one model the
first reward is independently $+1$ or $-1$ in either state. In the
other it is independently $+1$ or $-1$ in state zero and $+2$ or $-2$
in state one. Both have the same zero reward mean, transition and signal
kernel. At the next decision, choosing the state correctly pays one.
The optimal success probability is $4/5$ in the first model and one in
the second, because absolute reward reveals the state. Their full reward
observation laws differ; the example proves insufficiency of the compressed
mean-reward interface, not nonidentification from the full marked law.

For an exact two-state example take
\[
 E=\begin{pmatrix}4/5&1/5\\1/5&4/5\end{pmatrix},\quad
 \mu=\begin{pmatrix}3/5\\2/5\end{pmatrix},\quad
 T_0=\begin{pmatrix}4/5&1/5\\3/10&7/10\end{pmatrix},\quad
 T_1=\begin{pmatrix}3/5&2/5\\1/10&9/10\end{pmatrix}.
\]
Computing $J_0,J_1$ and applying the inverse recovers both transitions
exactly. If the two columns of $E$ instead coincide, observations
cannot distinguish the states through this construction. If the
columns become nearly equal, $L$ becomes large and finite-sample
errors are amplified twice in $LJ_aL'$. A fitted state dimension
that exceeds the measurement rank can therefore be computationally
feasible and scientifically unsupported.

For a logger using only the private state, the analogous joint matrix is
$E\operatorname{diag}\{\mu_t(s)q(a\mid s)\}T_aE'$ before
normalization. That selection factor cannot be ignored. Extra observed
proxies, a validated logging model or interventions can identify it
under additional restrictions; the action-conditioned matrix alone is
not the randomized matrix above.

\section{Proxy value identification without a complete latent fit}
\label{r3:sequential-proxies}
Past and current observations can sometimes serve as two views of the
hidden state. The matrix identification result of
\citet[Section 3, Theorem 1]{tennenholtz2020} uses specific
invertibility conditions and the declared behavior/evaluation information
sets. An independent hidden state at every date can destroy the needed
past-to-current rank: a long history is then not automatically an
informative proxy for today's private confounder.

The bridge approach makes the connection to
Chapter~\ref{r3:inverse-learning} more explicit.
\citet[Section 4, Assumptions 1--2 and Theorem 2]{bennettprox2023}
uses a hybrid law $P_t^*$ that applies the target policy before $t$
and the logging policy from $t$ onward. For a deterministic target,
let $D_t$ be the action that its observed information would select.
Its negative-control restriction has the form
\[
 (Z_t,A_t)\perp(W_t,D_t,Y_t(a))\mid U_t
 \quad\text{under }P_t^*,
\]
for the allowed future outcome variables $Y_t$. Here $Z_t$ is the
action-side proxy, $W_t$ the outcome-side proxy and $U_t$ the sufficient
hidden confounder, possibly augmented by observed context. This is
stronger than saying both proxies predict repayment. In particular,
giving the target policy extra history can invalidate independence of
$Z_t$ and $D_t$, even if the physical environment is unchanged.

The source's Assumption 2 requires existence of \emph{latent} action
and outcome bridges $q_t^\dagger,h_t^\dagger$ satisfying
\begin{align*}
 \E_t^*[q_t^\dagger(Z_t,A_t)\mid U_t,A_t=a]
       &=P_t^*(A_t=a\mid U_t)^{-1},\\
 \E_t^*[h_t^\dagger(W_t,A_t)\mid U_t,A_t=a]
       &=\E_t^*[\mathbf1(D_t=A_t)Y_t\mid U_t,A_t=a].
\end{align*}
This includes latent action positivity. It is stronger than merely finding
solutions of observed conditional equations, and is separate from the
negative-control independence. Under these two assumptions its Theorem 2
allows solutions $q_t,h_t$ of the corresponding observable equations
\begin{align*}
 \E_t^*[q_t(Z_t,A_t)\mid W_t,A_t=a]
       &=P_t^*(A_t=a\mid W_t)^{-1},\\
 \E_t^*[h_t(W_t,A_t)\mid Z_t,A_t=a]
       &=\E_t^*[\mathbf1(D_t=A_t)Y_t\mid Z_t,A_t=a].
\end{align*}
The outcome on the right is recursively constructed from later
bridges. These observed solutions need not be asserted to be the latent
bridges themselves. The theorem uses latent existence and the independence
restriction to justify the causal target of the observed solution.
The static augmented-bridge cancellation can then be applied
at each hybrid-law step. The resulting observed-trajectory functional
identifies policy value without assigning a uniquely interpretable
latent state to each account. Bridge existence, support and estimation
stability remain necessary; the identification result is not a claim
that every matrix inverse or neural minimax fit will be reliable.

In the source's tabular example with $U_t=S_t$, $Z_t=O_{t-1}$,
$W_t=O_t$ and a policy using only $O_t$, the outcome bridge requires
solving a linear system through $E_t'=P(O_t\mid S_t)'$.
The action bridge additionally involves the backward state kernel
$P_t^*(S_{t-1}\mid S_t)$ \citep[Proposition 1]{bennettprox2023}.
Full ranks provide a sufficient solvability route, but the backward
kernel depends on the induced policy law. It must be checked for that
law; invertibility of each action-specific transition should not be
substituted for checking a policy-dependent mixture without proof.
For example, $I_2$ and the two-state swap matrix are both invertible,
whereas their equal mixture has rank one.
This is the same distinction between structural identification and
inverse stability encountered in NPIV, now at each decision date.

\section{Estimating a controlled likelihood}
\label{r3:controlled-estimation}
For finite states, let $K_t(s,s')$ collect the action, transition,
reward and next-emission factors of Equation~\eqref{r3:controlled-law}.
The forward recursion multiplies a row distribution by $K_t$, normalizes
by $c_t$, and accumulates $\log L=\log c_{\rm initial}+\sum_t\log c_t$.
Scaling at each step avoids numerical underflow. A backward recursion
gives smoothed state probabilities $\gamma_t(s)$ and adjacent-state
probabilities $\xi_t(s,s')$.

An EM update maximizes the conditional expectation of the complete
log likelihood. When parameters separate and transitions are unrestricted
row probabilities, its action-specific transition update is
\begin{equation}
 \widehat T_a(s,s')=
 \frac{\sum_{i,t}\mathbf1(A_{it}=a)\xi_{it}(s,s')}
      {\sum_{i,t}\mathbf1(A_{it}=a)\gamma_{it}(s)}.
 \label{r3:controlled-em}
\end{equation}
Emission, reward and logger updates maximize their own expected log
factors when that factorization is valid. The sufficient counts generalize
the HMM forward--backward construction of \citet{rabiner1989}.
If the logger uses private state, omitting its likelihood factor changes
$\gamma,\xi$ and therefore every update, even if actions themselves
have no parameters of interest.

The usual likelihood ascent follows by Jensen's inequality applied to
the latent posterior at the old parameter: increasing its expected
complete log likelihood increases a lower bound tight at the old
parameter. Exact ascent does not guarantee the global maximum, a
correct number of states or causal identification. Multiple starts,
held-out predictive checks, nearly empty states, rank diagnostics,
boundary transitions and sensitivity to priors or penalties are material
to interpreting a fit. Bayesian estimation instead averages over a
parameter posterior proportional to $L(\theta,\eta)p(\theta,\eta)$;
it does not turn a likelihood equivalence into data identification.

\subsection{Continuous states and numerical alternatives}
For a linear Gaussian model with known/exogenous input,
\[
 S_{t+1}=FS_t+Ba_t+\eta_{t+1},\qquad O_{t+1}=CS_{t+1}+\nu_{t+1},
\]
Gaussian conditioning gives predictive mean $m^-=Fm_t+Ba_t$,
covariance $P^-=FP_tF'+Q$, innovation covariance
$V=CP^-C'+R$, gain $K=P^-C'V^{-1}$, and
\[
 m_{t+1}=m^-+K(O_{t+1}-Cm^-),\qquad
 P_{t+1}=P^--KVK'.
\]
The innovation density also supplies the likelihood increment.
This is the conditional-Gaussian derivation of
\citet[Theorem 4.2 and Theorem 12.1]{sarkka2013} with a deterministic
input added. State-dependent observation selection or a privately
chosen discrete action can make the conditional state law non-Gaussian,
even when the transition and emission equations are linear.

Nonlinear Gaussian filters approximate moments through linearization,
sigma points or quadrature. Particle methods approximate distributions
using weighted samples; variational methods optimize within a restricted
distribution family. These approximations have different failure modes.
A small covariance from a Gaussian approximation can miss a second
mode; a particle cloud can lose a low-probability state that matters
for a rare default; a variational objective can be accurate for its
loss and inaccurate for policy tail risk.

For sequential importance sampling, a particle transition proposal
$q^{\rm prop}(s'\mid s,\text{new data})$ requires the incremental
weight
\[
 w_{t+1}\ \propto\ w_t
 \frac{q_\eta(a_t\mid s,H_t)T_{a_t}(s,s')
        F(r_t\mid s,a_t,s')B(o_{t+1}\mid s',a_t)}
      {q^{\rm prop}(s'\mid s,\text{new data})}.
\]
It is the target/proposal path-density ratio
\citep[Section 7.3]{sarkka2013}, with the private logger retained.
For a known observed-history action probability, its factor is constant
across state particles and cancels in normalized filtering weights.
Under intervention it is replaced by the specified policy. State weights,
parameter uncertainty and policy uncertainty are separate objects.
Resampling effective sample size is a useful weight diagnostic, not
a bound on the error of a default-probability estimate.

\section{Errors in the model and errors in numerical solution}
\label{r3:value-errors}
An explicit bound helps decide where computational precision matters.
For a fixed Markov policy on a common finite state space,
$v=r+\gamma Pv$ with $0<\gamma<1$.
For an alternative $(\widetilde r,\widetilde P)$, suppose
$\|r-\widetilde r\|_\infty\le\epsilon_r$,
$\sup_s\|P(s,\cdot)-\widetilde P(s,\cdot)\|_1\le\epsilon_P$
and both rewards have magnitude at most $R$. Subtract the two Bellman
equations and use $\|\widetilde v\|_\infty\le R/(1-\gamma)$:
\begin{equation}
 \|v-\widetilde v\|_\infty
 \le \frac{\epsilon_r}{1-\gamma}
     +\frac{\gamma R\epsilon_P}{(1-\gamma)^2}.
 \label{r3:simulation-bound}
\end{equation}
An initial-distribution error contributes at most
$R\|\mu-\widetilde\mu\|_1/(1-\gamma)$.
This project bound assumes that the compared kernels represent the
same causal state/action system. It cannot align two causally
nonidentified models by fiat. A POMDP extension needs a valid comparison
of its induced belief processes and information sets.

For a contraction $\mathcal T$, an approximate value $\widehat v$
with Bellman residual at most $\epsilon_B$ satisfies
$\|\widehat v-v\|_\infty\le\epsilon_B/(1-\gamma)$: insert
$\mathcal T\widehat v$ between $\widehat v$ and $\mathcal Tv$
and rearrange. Numerical tolerance can therefore be chosen in the
units of the decision value. It should be smaller than the material
statistical and model uncertainty, but its smallness establishes
neither of them.

Dense forward/backward fitting costs $O(NTK^2)$ for $N$ independent
trajectories; declared sparse transitions can reduce the per-step
matrix cost. Trajectories can be partitioned and sufficient expected
counts combined. A long single series still carries temporal dependence,
and state sparsity is a model restriction if imposed for speed.
Numerical adapters should report the likelihood approximation, rank,
solver tolerances and resource measurements separately from causal
qualification.

The following comparisons put those obligations into concrete estimation
problems. Unknown emissions first require an identification argument for
the fitted model. Private logging then changes both the information needed
and the likelihood. An exact Gaussian construction supplies a numerical
reference before the discussion turns to customer choice. The final
servicing studies compare additional audit information and a nonlinear
model on explicitly restricted policy classes.



% SOURCE docs/release/2026-10-29/source/docs/writing-repair/2026-10-29/src/e764484b76fb2cf1.tex CHAR 0
\section{Unknown emissions: identification for the model actually fitted}
\label{s4:unknown-emission}
A known measurement matrix is a valuable baseline, but it does not justify
estimating that matrix from the same unvalidated records. A small controlled
model permits a constructive extension. Consider two states and two signals,
an unknown invertible emission matrix $E$ with columns summing to one,
unknown row-stochastic transitions $T_a$, and an unknown positive initial
law $\mu_0$. Require one invertible baseline transition $T_0$. Across
independent trajectories, actions are assigned independently of states,
signals and previous innovations, with positive probabilities. Signal
draws are conditionally independent given their states and invariant across
actions. Conditional on the current state/action, a binary repayment reward
is independent of signal and next-state draws. These restrictions describe
one initial supported POMDP estimation experiment; they are not assumptions
learned from EM convergence.

Condition on $A_0=A_1=0$. Independent assignment makes this the observation
law under two imposed baseline actions. Conditional on $S_1$, the three
signals $O_0,O_1,O_2$ are independent. Let
\[
 \mu_1=T_0'\mu_0,\qquad
 L=E\operatorname{diag}(\mu_0)T_0\operatorname{diag}(\mu_1)^{-1},
 \qquad R=ET_0'.
\]
The columns of $L$ and $R$ are the conditional past/future signal laws
given $S_1$; both are invertible. For each middle signal $j$, define the
observable matrix $C_j(i,k)=P(O_0=i,O_1=j,O_2=k\mid A_0=A_1=0)$.
Then
\begin{align}
 P_{02}:=\sum_jC_j&=L\operatorname{diag}(\mu_1)R',\nonumber\\
 C_jP_{02}^{-1}&=L\operatorname{diag}(E_{j1},E_{j2})L^{-1}.
 \label{s4:controlled-spectral}
\end{align}
Invertibility of a binary emission matrix means the two eigenvalues for
either informative row differ. Its eigenvectors recover $L$ up to
column scaling/permutation. Normalize each column to sum to one; recover
each row of $E$ as the diagonal of $L^{-1}C_jP_{02}^{-1}L$.
The same permutation applies to every row. Now
$\mu_0=E^{-1}p_{O_0}$ and, for each supported action,
\begin{equation}
 T_a=\operatorname{diag}(\mu_0)^{-1}E^{-1}
 P(O_0,O_1\mid A_0=a)E^{-\prime}.
 \label{s4:unknown-transition-recovery}
\end{equation}
For $b_a(s)=P(R_0=1\mid S_0=s,A_0=a)$, the vector
$k_a(o)=P(O_0=o,R_0=1\mid A_0=a)$ gives
$b_a=\operatorname{diag}(\mu_0)^{-1}E^{-1}k_a$.
Thus the unknown-emission, transition and binary-reward model used below
is identified, up to one common state relabeling. No stationary initial
law was required. The state label can be fixed by ordering distress-signal
probabilities; the monetary policy value itself is permutation invariant.

This proof is a project construction using three observed views.
\citet[Sections 3.1--3.2]{spectral2016} develops a broader controlled
spectral route with action-conditioned views, rank assumptions, a
memoryless logging policy and single-trajectory mixing bounds. Its small
singular values enter the finite-sample requirements explicitly. That is
the relevant comparison to an uncontrolled HMM theorem: the action's
dependence on observations changes the views. Augmenting the state by
the action and asserting generic HMM full rank can instead introduce a
rank-deficient transition matrix.

Three limitations follow from the proof itself. An almost singular $E$,
$T_0$ or $P_{02}$ amplifies sampling error; identification is not a
conditioning guarantee. With more hidden states than signal categories,
this two-view inversion fails, even though a different multi-period theorem
may apply. Finally, conditioning on baseline actions chosen using private
information need not leave the experimental law used in the first line.
A logger-specific joint model or additional experimental/proxy information
is then necessary. Under the special private-state-only logger, known
emissions can sometimes recover $\mu(s)q(a\mid s)$ from the joint matrix
and recover $T_a$ by row normalization. Failure of the naive conditional
likelihood is therefore not a universal impossibility theorem for every
private-logger model.

\section{From an identified law to fitted parameters and value uncertainty}
\label{s4:latent-inference}
In the binary experiment write $\vartheta\in\mathbb R^{11}$ for logits
of the initial state probability, two emission probabilities, four
action-specific transition probabilities and four repayment probabilities.
The data for trajectory $i$ are $o_{i,0:H-1},a_{i,0:H-1},r_{i,0:H-1}$.
Conditional on randomized actions, its likelihood is
\[
 L_i(\vartheta)=\sum_{s_{0:H-1}}\mu(s_0)
 \prod_{t=0}^{H-1}E_{o_t,s_t}
     b_{a_t}(s_t)^{r_t}\{1-b_{a_t}(s_t)\}^{1-r_t}
 \prod_{t=0}^{H-2}T_{a_t}(s_t,s_{t+1}).
\]
All final-date reward and emission terms remain present; no final-date
transition is observed in this experiment. The forward recursion sums
this likelihood in $O(HK^2)$ per trajectory and uses a separate scaling
constant at every date. Backward recursion gives $\gamma_{it}(s)$ and
$\xi_{it}(s,s')$ as above. Their expected counts update each Bernoulli
probability. A probability box $[0.002,0.998]$ gives the exact constrained
complete-data maximizer by clipping the expected-count ratio.

Three starts of scaled EM supply initial fits; a direct L-BFGS-B
likelihood maximization using an analytic score provides a numerical
challenger and refinement. For a transition logit $\vartheta_{as}$, the
trajectory score follows by differentiating the hidden-path sum:
\[
 \partial_{\vartheta_{as}}\log L_i
 =\sum_{t=0}^{H-2}\ind(a_{it}=a)
       \{\xi_{it}(s,1)-T_a(s,1)\gamma_{it}(s)\}.
\]
Emission, initial-law and reward logits give analogous posterior success
counts minus fitted probability times posterior trial counts. This is
also a direct derivation of the conditional complete-data score identity.
Finite differencing this analytic score gives the observed information,
including the dependence of the posterior weights on parameters. Treating
the smoothed states as known would omit that missing-state uncertainty.

The implementation uses a centered logit step of $10^{-4}$ for the
information and value gradient. A separate check differentiates the direct
observed likelihood at step $10^{-5}$ and compares with the posterior-count
score; a hidden-path sum independently checks the forward likelihood.
The finite-difference truncation and rounding errors are numerical errors,
not sampling variation. Increasing state dimension or approaching singular
information requires renewed step-size/conditioning analysis. The reported
local curvature diagnostics do not certify a global optimum.

\subsection{The regular sampling argument and its failures}
Take $N$ independent complete trajectories and fixed $H$. Suppose the true
probabilities are in the interior of the declared compact box, rank and
emission separation hold, labels are fixed, and the information matrix
$I=\E[s_i s_i']=-\E[\partial_\vartheta s_i]$ is nonsingular.
For finite alphabets there are finitely many observed trajectory patterns.
Their empirical frequencies obey a uniform law of large numbers; positivity
bounds the pattern log likelihoods and derivatives. The expected log
likelihood is uniquely maximized at the identified law by nonnegativity
of Kullback--Leibler divergence. A global approximate maximizer is
therefore consistent on the label-ordered compact parameter space.
Taylor expansion of its score yields
\begin{equation}
 \sqrt N(\widehat\vartheta-\vartheta_0)
 =I^{-1}\frac1{\sqrt N}\sum_i s_i(\vartheta_0)+o_p(1).
 \label{s4:latent-mle-clt}
\end{equation}
The iid finite-pattern CLT and the information identity then give
covariance $I^{-1}$. This supplies an argument for the actual sampling
law, not a generic claim about a single dependent series or growing $H$.
Multiple EM starts do not prove global maximization. Boundary solutions,
singular information, optimizer caps and disagreement between starts must
remain reported failures or conditional fits.

For a policy $\pi$ fixed before evaluation, let $V_\pi(\vartheta)$ be
the expected discounted recorded financial reward. For finite alphabets
and fixed horizon, summing over hidden paths and observed histories is a
finite smooth function of the environment probabilities, even if $\pi$
uses a threshold. Its history-to-action map is held fixed. Consequently
\[
 \widehat{\operatorname{Var}}(\widehat V_{\pi_1}-\widehat V_{\pi_0})
 =N^{-1}\nabla(V_{\pi_1}-V_{\pi_0})'
       \widehat I^{-1}\nabla(V_{\pi_1}-V_{\pi_0}).
\]
The contrast gradient retains covariance between the policies. Resampling
whole trajectories and refitting all environment parameters is a bootstrap
alternative in this regular regime. It keeps the policy fixed and repeats
initialization, optimization and label handling. A bootstrap that discards
capped fits or recomputes an optimal policy answers a different question.

In the implemented example, $\pi_1$ assists exactly when the binary signal
indicates distress, while $\pi_0$ always uses action zero. Each repayment
pays 100 currency units and assistance costs 12; $H=8$, $\gamma=0.99$.
For a reactive policy, define
\[
 P_\pi(s,s')=\sum_o E_{os}T_{\pi(o)}(s,s'),\qquad
 r_\pi(s)=\sum_o E_{os}\{100b_{\pi(o)}(s)-12\pi(o)\}.
\]
The exact finite-horizon recursion $v_h=r_\pi+\gamma P_\pi v_{h-1}$,
$v_0=0$, gives $V_\pi=\mu'v_H$. This is a fixed observable policy
implemented on a partially observed environment; it does not claim to
be the optimal belief policy. A policy learned from a fitted filter must
retain that trained filter's history-to-action map during all subsequent
environment refits, or use an explicit separate policy-selection design.

\subsection{Numerical tolerance in the same units as policy value}
Within a basin where the negative average log-likelihood Hessian has
smallest eigenvalue at least $\kappa>0$, an average score residual
$\epsilon_s$ bounds parameter optimization error by
$\epsilon_s/\kappa$. A gradient bound $\|\nabla V\|\le G_V$ then
gives value error at most $G_V\epsilon_s/\kappa$. For asymptotic equivalence
in Equation~\eqref{s4:latent-mle-clt}, require
$\epsilon_s=o(N^{-1/2})$ as well as selection of the correct basin.
The observed eigenvalue is a local diagnostic, not proof of a uniform
population curvature bound. Solver success and small likelihood increments
alone do not establish an acceptable monetary error.

Spectral moment inversion is a meaningful alternative: it has a constructive
population solution and avoids iterative latent-path fitting, but inverses
and eigengaps can amplify noisy empirical moments into negative probabilities.
Constrained likelihood uses all recorded patterns and enforces valid
probabilities at a cost in nonconvex optimization. Spectral initialization,
likelihood refinement and the exact small-state reference therefore serve
different purposes. Gaussian filtering, particle likelihood and variational
inference retain the approximation distinctions in the preceding section;
none inherits this finite-alphabet likelihood theorem without a new argument.

\section{A value-error bound on the actual partially observed process}
\label{s4:pomdp-error}
The fully observed bound in Equation~\eqref{r3:simulation-bound} does not
by itself compare two belief processes. A path-law argument does. Align
the same latent labels and observed alphabet in two controlled models.
Let the initial joint state/signal laws differ by at most $\delta_0$ in
$L_1$. Let the joint one-step kernel for reward, next state and next signal,
given current state and action, differ by at most $\delta_K$ uniformly.
Use the \emph{same} policy kernel on every observed history in both models.
Its state in this argument is the complete history together with the latent
state, not just the current observed signal.

Markov kernels contract $L_1$ distance. Insert the first model's one-step
kernel between the two propagated path laws and apply the triangle
inequality. Induction bounds the path-law distance through reward $t$ by
$\delta_0+(t+1)\delta_K$. If $|R_t|\le R$, integration gives
\begin{equation}
 |V_\pi-\widetilde V_\pi|
 \le R\sum_{t=0}^{H-1}\gamma^t
               \{\delta_0+(t+1)\delta_K\}
 \le \frac{R\delta_0}{1-\gamma}
       +\frac{R\delta_K}{(1-\gamma)^2}.
 \label{s4:pomdp-path-bound}
\end{equation}
For factored transition, reward and emission kernels, telescoping the
product measures bounds $\delta_K$ by the sum of their uniform $L_1$
errors. The bound is often conservative, but it explicitly respects partial
observation and policy information. Comparing different optimal policies
or updating the stored policy with each model does not satisfy its premise.

For numerical evaluation of one fixed model/policy, let the backward value
recursion at time $t$ have sup-norm residual $\epsilon_t$. Applying the
same contraction argument backward bounds initial value error by
$\sum_t\gamma^t\epsilon_t$. An infinite-horizon discounted belief
operator gives $\epsilon/(1-\gamma)$ for a uniform Bellman residual.
Both require the operator to be the one induced by the specified observation
and policy process. A grid-spacing comparison alone is not a certified
residual bound. Model error, sampling error, policy selection and numerical
error are distinct terms in a financial decision report.


% SOURCE docs/release/2026-10-29/source/docs/writing-repair/2026-10-29/src/ce2741b43a3eca45.tex CHAR 0
\section{Private decisions and independently validated measurements}
\label{s5:private-logger}
An arrears officer may see information absent from the warehouse. Adding
states to a predictive HMM does not explain that information. A tractable
alternative obtains an independent labelled measurement sample and models
the action mechanism explicitly. Let $S_t$ have $K$ states, $O_t$ have $M$
signal categories, and $A_t$ have two actions. On one decision clock assume
\[
 O_t\mid S_t\sim E_{\cdot S_t},\quad
 A_t\mid S_t\sim q_{S_t},\quad
 R_t\mid S_t,A_t\sim N(r_{A_t,S_t},\sigma_R^2),\quad
 S_{t+1}\mid S_t,A_t\sim T_{A_t}(S_t,\cdot).
\]
The displayed factors assert conditional independence of signal, logger
randomness and reward noise given the indicated parents. Reward occurs before
the next signal. Conditional independence, time invariance, state labels and
the action/transition modularity are scientific premises. The validation
sample observes labelled $(S,O)$ under the same emission law; labels sampled
uniformly across states need not reproduce the population state frequencies.
Their conditional emission likelihood is the relevant contribution.

\subsection{Identification from unnormalized marked moments}
Let $\mu_s=P(S_t=s)$ and define the observable matrix and reward vector
\[
 J_a(o,o')=P(O_t=o,A_t=a,O_{t+1}=o'),\qquad
 k_a(o)=\E[R_t\mathbf1\{O_t=o,A_t=a\}].
\]
The factorization implies
\[
 J_a=E\operatorname{diag}(\mu\odot q(a))T_aE',\qquad
 k_a=E\operatorname{diag}(\mu\odot q(a))r_a.
\]
If the validated $M\times K$ emission matrix has full column rank, let
$L=(E'E)^{-1}E'$. Then
\begin{equation}
 B_a=LJ_aL'=\operatorname{diag}(\mu\odot q(a))T_a,
 \quad w_a=B_a\mathbf1,\quad
 \mu=\sum_a w_a,\quad q_s(a)=w_{as}/\mu_s.
 \label{s5:marked-identification}
\end{equation}
Positive $w_{as}$ permits recovery of $T_a$ by row division and
$r_a=\operatorname{diag}(w_a)^{-1}Lk_a$. This is a constructive sufficient
identification proof using population moments. Conditioning $J_a$ on $A=a$
without retaining $P(A=a)$ discards the action-specific state mass required
by the inversion. An unvalidated emission law may be identified by other
longer-history restrictions; failure of this particular route is not a
universal latent-identification impossibility theorem.

The smallest singular value of $E$ and the minimum state-action mass control
error amplification. Perturbing $J_a$ with $E$ fixed gives
$\|dB_a\|\le\|L\|^2\|dJ_a\|$; perturbing $E$ changes both inverses.
Rank diagnostics can reveal weak separation, but a labelled audit from a
different measurement process may give a precise, wrong $E$. Validation
invariance and label meaning therefore remain explicit restrictions.
More data from a biased validation source do not remove that bias.

\subsection{Joint likelihood, score and two-sample uncertainty}
For one observed trajectory the likelihood is
\[
 L_i(\vartheta)=\sum_{s_{0:H-1}}\mu_{s_0}
 \prod_{t=0}^{H-1}E_{o_t s_t}q_{s_t}(a_t)
       f_R(r_t\mid s_t,a_t)
 \prod_{t=0}^{H-2}T_{a_t}(s_t,s_{t+1}).
\]
The joint log likelihood adds $\sum_{j=1}^{M_v}\log E_{o_j^v s_j^v}$
from the independent validation sample. Softmax row coordinates enforce
probability constraints; Gaussian reward means are scaled by their known
standard deviation for numerical conditioning. Forward probabilities are
normalized at every step. Reward log-density maxima are subtracted and
added back to the log likelihood, so the stabilized computation evaluates
the same function.

Differentiating the finite hidden-path sum gives
$\nabla\log L_i=\E_\vartheta[\nabla\log L_i^{\rm complete}\mid O,A,R]$.
For transition-logit coordinate $(a,s,j)$, for example, the score is
\[
 \sum_t\mathbf1\{A_t=a\}
 \{\xi_t(s,j)-\gamma_t(s)T_a(s,j)\},
\]
where $\gamma,\xi$ are state and adjacent-state smoothing probabilities.
Emission and logger coordinates have the same expected-count minus
expected-total-times-probability form. Reward-mean scores are posterior
weighted residuals divided by $\sigma_R^2$. The validation emission score
uses the observed state label, not a smoothing probability. Exhaustive
hidden-path enumeration checks the likelihood and a separate derivative
checks this analytic score to below $10^{-9}$ in the small reference.

Let trajectory and validation scores be $s_i$ and $u_j$. With fixed horizon,
an identified interior parameter, differentiability, finite moments and
uniform local information, Taylor expansion of the total score gives
\[
 \hat\vartheta-\vartheta_0=
 A^{-1}\left(\sum_{i=1}^Ns_i+\sum_{j=1}^{M_v}u_j\right)
 +o_p(N^{-1/2}),\qquad
 A=N I_o+M_v I_v,
\]
when $M_v/N$ converges to a positive finite constant. Under the correctly
specified model the covariance is $A^{-1}$. A two-sample sandwich uses
$A^{-1}(N J_o+M_v J_v)A^{-1}$, centering each independent score sample
separately. It estimates local sampling uncertainty under possible
likelihood misspecification; it does not restore a causal interpretation
when the model factorization is wrong. If validation alone identifies a
direction, letting $M_v/N\to0$ may change that direction's rate.

Partitioning information into structural and emission parameters shows why
plugging in calibration as known can understate variance. The effective
structural information is the Schur complement
$A_{\beta\beta}-A_{\beta E}A_{EE}^{-1}A_{E\beta}$, not
$A_{\beta\beta}$. For a fixed signal policy $\pi(a\mid o)$,
$\bar\pi(a\mid s)=\sum_oE_{os}\pi(a\mid o)$ determines the target
transition and reward. The value gradient therefore includes the emission
derivative as well as transition, reward and initial-state derivatives.
The historical logger disappears from the intervention recursion but remains
in the likelihood that identifies those parameters.

\subsection{Numerical and inferential challengers}
The new finite experiment fits three states, four signals, private logging,
eight decisions, 600 customer trajectories and either 90 or 600 validation
labels. It estimates 32 parameters using two data-based starts. Reward noise
standard deviation is known at 12. The true fixed observable-policy contrast
is $59.5645$. With 90 labels, bias is $2.0798$ and coverage .850; with 600,
bias is $.1926$ and coverage .875. These denominators include all 80 requested
replications per case, including four and five boundary/numerical failures.
Those failures yield unavailable inference. Model, sandwich and
fixed-emission approximations happen to give the same coverage here; the
reward is informative about state and the measured calibration contribution
is small. This experiment does not demonstrate a large calibration variance
penalty. Omitting the private logger on the first 20 datasets yields biases
of $-36.51$ and $-41.12$, respectively.

The first 600-label dataset was selected before a 199-refit parametric
bootstrap challenger. It resimulates complete trajectories and independent
emission measurements conditional on the audited label counts, then refits
the complete model. Seventeen refits hit a boundary or failed, exceeding the
prespecified 5\% availability limit, so no percentile/basic interval is
released. The standard deviation among 182 usable refits is $10.20$, only a
descriptive conditional diagnostic. Extra starts change the fitted contrast
by $.0012$ and the average negative log likelihood by less than $10^{-10}$.
Information differencing steps $3\times10^{-5},10^{-4},3\times10^{-4}$
change the model standard error by less than $4\times10^{-8}$. The local
information condition number is about 1,929. These checks make numerical
roundoff an implausible sole explanation of the observed undercoverage;
they do not prove the global optimum or finite-sample normality.

The method decision is therefore explicit: this validated finite route is
an identified research fit with a useful value functional, but its current
normal/bootstrapped interval is not a routine calibrated product default.
The larger-sample and audited-trajectory experiments in
Section~\ref{s6:servicing} provide selected alternatives; they do not
retroactively qualify this 600-trajectory interval. Weak-state/profile
inference and bank measurement validation remain distinct requirements. The observed
failure is part of the proposal's method hierarchy, not a hidden exception.

\subsection{Proxies identify a different object with different information}
Sequential proxy bridges can identify value without recovering a uniquely
labelled latent transition model \citep{bennettprox2023,zhangprox2024}.
Their conditional independence and completeness restrictions are different
from labelled measurement validation. The two-stage dynamic-regime source
identifies joint counterfactual quantities along compatible treatment
histories under its stated bridge restrictions; it does not identify every
cross-world joint law. Bridge existence must be asserted or proved separately
from uniqueness. For an operator $K$, injectivity means $Kh=0$ implies
$h=0$; it does not mean that every observed right-hand side lies in
$\operatorname{range}(K)$.

This matters for missing rewards. The 2026 missingness-aware MDP preprint
uses augmented state, a shadow next-state variable, exclusion and completeness
to motivate a bridge \citep[Assumptions 3.1--3.4, Theorem 3.6]{wei2026}.
Its general existence discussion cannot be adopted from injectivity alone.
The operator contract therefore needs an integrable bridge solution and
the relevant range condition before uniqueness can justify identification.
An inverse-problem regularization analysis is then needed to control
estimation error. In a finite square full-rank system
existence follows for every right-hand side; an infinite-dimensional
conditional expectation operator needs more. Overlap, residual and rank
diagnostics are conditional checks. They cannot certify the unobserved
exclusion and completeness assumptions that make an MNAR query causal.


% SOURCE docs/release/2026-10-29/source/docs/writing-repair/2026-10-29/src/c9cec5029cff1de4.tex CHAR 0
\section{An exactly Gaussian controlled model with private logging}
\label{s5:continuous-controlled}
The finite-state construction leaves a numerical question: can a controlled
continuous-state likelihood retain the private action mechanism and still
admit exact filtering and derivatives? A scalar Gaussian reference supplies
an affirmative answer under specific measurement restrictions. Let
\begin{align}
 X_0&\sim N(0,1), & O_t&=X_t+\nu_t, & A_t&=kX_t+\eta_t,\\
 R_t&=cX_t+dA_t+\epsilon_t, &
 X_{t+1}&=fX_t+bA_t+\omega_{t+1}.
 \label{s5:gaussian-private-law}
\end{align}
All primitive noises are independent with known positive variances
$(r_o,r_a,r_r,q)$. The unknown coefficients are $(f,b,k,c,d)$.
The signal loading is calibrated to one and the initial law fixes the state
scale. Actions are continuous here, such as a scaled assistance amount.
Known noise variances and invariant measurements are a substantial price for
this exact reference, not defaults for arbitrary bank feeds.

\subsection{Identification and the likelihood actually evaluated}
At the initial date, $\operatorname{Cov}(O,A)=k$,
$\operatorname{Cov}(O,R)=c+dk$ and
$\operatorname{Cov}(A,R)=kc+d(k^2+r_a)$. Hence
\[
 d=\frac{\operatorname{Cov}(A,R)-k\operatorname{Cov}(O,R)}{r_a},
 \qquad c=\operatorname{Cov}(O,R)-dk.
\]
Similarly, writing $g_1=\operatorname{Cov}(O_1,O_0)=f+bk$ and
$g_2=\operatorname{Cov}(O_1,A_0)=fk+b(k^2+r_a)$ gives
$b=(g_2-kg_1)/r_a$ and $f=g_1-bk$. Independent customer trajectories
identify these moments. Positive private-action noise $r_a$ supplies the
independent variation separating action response from hidden state. If it
vanishes, these inversions fail. Shared action/reward noise or an unknown
measurement scale invalidates this particular argument. These project
identities establish the identification experiment used by the fit.

Given the predictive state law $N(m_t,P_t)$, define
\[
 z_t=(O_t,A_t,R_t-dA_t)',\quad C=(1,k,c)',\quad
 v_t=z_t-Cm_t,\quad S_t=P_tCC'+\operatorname{diag}(r_o,r_a,r_r).
\]
The transformation from $(O,A,R)$ to $z$ has determinant one, including
when $d$ is unknown. The innovation density therefore gives the observed
joint likelihood of signal, private action and reward:
\[
 \ell_t=-\tfrac12\{3\log(2\pi)+\log|S_t|+v_t'S_t^{-1}v_t\}.
\]
Gaussian conditioning yields $K_t=P_tC'S_t^{-1}$,
$m_t^+=m_t+K_tv_t$ and $P_t^+=P_t-K_tCP_t$. Propagate using the observed
action: $m_{t+1}=fm_t^++bA_t$ and $P_{t+1}=f^2P_t^++q$.
This extends the Gaussian conditioning/innovation likelihood of
\citet[Theorems 4.2 and 12.1]{sarkka2013} by explicitly including the
private logger as a measurement channel. Treating $A_t$ only as an exogenous
input omits information and changes the likelihood.

For parameter coordinate $j$, use a dot for differentiation and
$w_t=S_t^{-1}v_t$. Then
\begin{align}
 \dot v_t&=\dot z_t-\dot C m_t-C\dot m_t, &
 \dot S_t&=\dot P_tCC'+P_t(\dot CC'+C\dot C'),\\
 \dot\ell_t&=-\tfrac12\operatorname{tr}(S_t^{-1}\dot S_t)
       +\tfrac12w_t'\dot S_tw_t-\dot v_t'w_t,\\
 \dot K_t&=(\dot P_tC'+P_t\dot C'-K_t\dot S_t)S_t^{-1},\\
 \dot m_t^+&=\dot m_t+\dot K_tv_t+K_t\dot v_t, &
 \dot P_t^+&=\dot P_t-\dot K_tCP_t-K_t\dot CP_t-K_tC\dot P_t,\\
 \dot m_{t+1}&=f\dot m_t^++\dot f m_t^++\dot b A_t, &
 \dot P_{t+1}&=f^2\dot P_t^++2f\dot f P_t^+.
 \label{s5:gaussian-score}
\end{align}
The derivative of the reward transform is
$\dot z_t=(0,0,-A_t\dot d)'$; omitting it differentiates the wrong
likelihood. Initial sensitivities are zero because the initial law is
fixed. Unknown noise or initial parameters would add their own terms.
The implementation uses solves for innovation systems, accumulates the
analytic score, and finite-differences that score for observed information.
The covariance update is symmetrized in matrix generalizations; a square-root
factorization is the preferred extension for poorly conditioned systems.

An independent check constructs all observed variables as linear combinations
of the initial state and primitive noises, forms their full joint covariance,
and evaluates a multivariate Gaussian density. This agrees with the sequential
likelihood; central differences independently check the analytic score.
These checks include the action channel and reward transform, not only a
Kalman recursion with fixed inputs. They meet the numerical obligation to
differentiate the likelihood being optimized.

\subsection{Intervention value and a complete fitted comparison}
Replace the historical logger by the fixed observable policy
$A_t=\alpha+\beta O_t$. Under this intervention
\[
 \bar x_{t+1}=(f+b\beta)\bar x_t+b\alpha,\qquad
 \bar r_t=(c+d\beta)\bar x_t+d\alpha,
 \quad V=\sum_{t=0}^{H-1}\gamma^t\bar r_t.
\]
The historical coefficient $k$ disappears from the intervention law.
Policy signal noise affects the distribution even though its mean is zero;
the displayed recursion is sufficient for a linear expected reward, not a
default-tail probability or nonlinear utility. Every prescribed action lies
in the historical Gaussian law's topological support. Nevertheless the
deterministic conditional policy is singular relative to its continuous
logging density, so ordinary importance weighting does not apply. This
identification rests on the structural functional form and identified
coefficients, not a bounded-density-ratio theorem. A finite
horizon needs no infinite-horizon stability; an infinite value would need
appropriate growth/discount conditions.

The experiment uses $H=12$, $N=800$, coefficients
$(.6,.35,.8,1.4,-.2)$ and noise variances $(q,r_o,r_a,r_r)=(.3,.4,.5,.7)$.
The policy has $(\alpha,\beta,\gamma)=(.4,-.2,.97)$, giving true value
$2.6785$ relative to the zero-action, zero-mean reward reference. Across
160 repetitions the joint Gaussian fit has bias $-.0020$ and nominal 95\%
coverage .9625, with no unusable fit. The Monte Carlo standard error near
this coverage is .015. The exogenous-input challenger omitting private
logging has value bias $.9531$. Only a point comparison is made for that
misspecified likelihood. The initial likelihood and gradient identity checks
are retained with the experiment, as are optimizer residuals and information
conditioning.

This comparison supplies a selected continuous route, not a blanket POMDP
estimator. For discrete private actions with a logistic state-dependent
logger, the posterior is generally non-Gaussian; the exact recursion above
no longer applies. An extended/unscented filter approximates moments and
may bias likelihood derivatives. Quadrature is a deterministic challenger
in low dimension, requiring node/tolerance sensitivity. Particle likelihood
keeps the logger in the weights, with independent particle replicates and
particle-number studies separating Monte Carlo variation from parameter
sampling uncertainty. Differentiating a resampling implementation requires
its own argument; a formal derivative of smooth transition densities is
not automatically the derivative of the stochastic objective actually used.

For a locally strongly curved average negative log likelihood with Hessian
lower bound $\kappa$, an average score error $\varepsilon_s$ moves the local
solution by at most approximately $\varepsilon_s/\kappa$. A value-gradient
bound $G$ then contributes $G\varepsilon_s/\kappa$ to numerical value error.
Likelihood approximation adds its own score error, and value quadrature or
simulation adds a separate term. This conditional error budget explains why
reducing solver tolerance cannot repair measurement misspecification. The
exact Gaussian route is a reusable benchmark for numerical adapters; richer
state models require renewed identification, derivative and uncertainty
qualification.


% SOURCE docs/release/2026-10-29/source/docs/writing-repair/2026-10-29/src/73ae9588bcea9fe8.tex CHAR 17460


\section{Customer choice is an inverse decision problem}
\label{r3:customer-choice}
The bank can optimize a supplied reward while a customer chooses
whether to refinance, accept an offer or default. Estimating that
customer's behavior reverses the direction of the decision problem:
choices and transitions are observed, and preferences are unknown.
Let $x$ be the customer's observed Markov state, $d$ a customer
action, $u_d(x)$ its flow utility, and $P_d$ its transition matrix.
Assume independent choice shocks with a specified type-I extreme-value
law, normalized so that the expected maximum has no additive constant,
and a specified discount $0\le\beta<1$. Then
\begin{align}
 v_d(x)&=u_d(x)+\beta(P_dV)(x),\nonumber\\
 V(x)&=\log\sum_d\exp v_d(x),\qquad
 p_d(x)=\frac{\exp v_d(x)}{\sum_j\exp v_j(x)}.
 \label{r3:ddc-logit}
\end{align}
The log odds identify $v_d-v_0=\log(p_d/p_0)$ under this model,
not the flow-utility difference alone. Future transitions enter the
choice-specific values. The conditional independence restriction on
next-period states and shocks, information available to customers,
discount and shock normalization are structural assumptions
\citep[Section 2]{kalouptsidi2021}.

Using $v_0=V+\log p_0$ gives
\begin{align}
 V&=(I-\beta P_0)^{-1}(u_0-\log p_0),\nonumber\\
 u_d&=(I-\beta P_d)(I-\beta P_0)^{-1}(u_0-\log p_0)+\log p_d.
 \label{r3:ddc-inversion}
\end{align}
The inverse exists because the Neumann series
$\sum_{k\ge0}\beta^kP_0^k$ converges. Thus observed CCPs and
transitions leave the vector $u_0$ free unless further restrictions
identify it. Equation~\eqref{r3:ddc-inversion} reconstructs the
reference-payoff argument in equations (3)--(5) of the cited paper.
A constant utility location and scale are familiar normalizations;
fixing an entire state-dependent payoff vector is more consequential.

\subsection{The same choices can imply different policy responses}
For any bounded state function $c$, define
\begin{equation}
 u_d^c=u_d+c-\beta P_dc,\qquad V^c=V+c.
 \label{r3:utility-equivalence}
\end{equation}
Then $v_d^c=v_d+c$ for every action at that state, so all choice
probabilities are unchanged. This is a constructive observational
equivalence, with a telescoping discounted reward interpretation.
If a counterfactual changes $P_d$ while retaining the two candidate
flow-utility models, the cancellation need not persist. It is precisely
the requested counterfactual that must be identified, rather than every
payoff or merely the historical choice probabilities
\citep{kalouptsidi2021}.

In a two-state example, let $\beta=1/2$, both historical action
transitions equal $I$, and all baseline utilities equal zero. Customers
choose each action with probability $1/2$. Taking $c=(0,2)'$ gives
equivalent utilities $u_d^c=(0,1)'$ for both actions. Now let action
1 move either state to state 1 while action 0 preserves the state.
The zero-utility model still predicts probability $1/2$ of action 1.
In the second model, state 1 has
$V_1=2(1+\log2)$ and, writing $z=\exp(V_0/2)$,
$z^2=z+2e$. Thus
\[
 p_1(0)=1-\frac{1}{z},\qquad
 z=\frac{1+\sqrt{1+8e}}2,
\]
which exceeds $1/2$. Both models fit the same historical choices and
transitions. The new transition makes their policy predictions differ.
In mortgage applications, a modification that changes repayment/default
transitions needs more identification than a well-fitted historical
refinancing-choice likelihood.

\section{Estimation alternatives and the bank--customer interface}
\label{r3:choice-estimation}
A full-solution likelihood estimates utility/transition parameters and
solves Equation~\eqref{r3:ddc-logit} at each parameter value. A CCP
route estimates choice probabilities and transitions first, then uses
an inversion such as Equation~\eqref{r3:ddc-inversion} with identifying
restrictions on payoffs. It trades repeated dynamic programming for
first-stage approximation and support requirements: very small CCPs
make log odds unstable, and unobserved customer types cannot be ignored
by pooling their nonlinear choice probabilities.

Bayesian estimation can integrate preference uncertainty, but the
likelihood still depends on a value solution. \citet[Sections 2.2--2.3]{imai2009}
compares conventional nested posterior simulation with an algorithm
that coordinates parameter simulation and approximate value updates.
Its convergence result has explicit regularity, compactness, memory
and bandwidth conditions (Assumptions 1--8 and Theorems 1--2).
Performing one approximate Bellman update inside arbitrary MCMC does
not inherit that result.

For the normalized logit model, a value error at most $\epsilon_V$
changes each $v_d$ by at most $\beta\epsilon_V$. Since log-sum-exp
is one-Lipschitz in sup norm, each log choice probability changes by
at most $2\beta\epsilon_V$. Combined with the Bellman-residual bound,
this gives a transparent likelihood error budget. It is a deterministic
approximation bound, not a sampling theorem for the fitted preferences.

The bank's action $a$ should remain distinct from the customer's
response $d$. An offer can change a contractual cost in $u_d(x;a)$,
the transition law $P_d(\cdot\mid x;a)$, or the information available
to the customer. A structural intervention must state which mechanism
changes and what remains invariant. Competitor repricing, funding
conditions and capacity can then alter the environment again. A
single-customer policy response conditional on those quantities is
not automatically a banking-group equilibrium effect. The product
must return the corresponding target and require an explicit equilibrium
model for the latter question.



% SOURCE docs/release/2026-10-29/source/docs/writing-repair/2026-10-29/src/bc194ea4c2c64949.tex CHAR 0
\section{A resolved information and inference design for servicing}
\label{s6:servicing}
The privately logged three-state model is useful precisely because it can
fail despite an identified population value. Resolving that failure requires
separating the statistical experiment from the optimizer and from the
proposed decision. Here the servicing decision is whether to offer temporary
assistance; the objective is discounted collected cash net of assistance and
operating cost. Signals arrive before the decision, while the cash rewards
used in this four-decision policy are settled after its decision horizon.
The policy may use signals and its own previous actions. It may not use the
later audited state or settled cash as contemporaneous information.

\subsection{What the retained failure actually shows}
For the original 600-trajectory experiment, all nine unavailable fits were
on a constrained transition-logit boundary. Reconstructing those exact seeds
reproduces the stored values, rather than diagnosing a different simulation.
With 90 emission labels, 76 of 80 fits were available: conditional coverage
was .8947, empirical value standard deviation 10.668 and mean model standard
error 9.657. With 600 labels, 75 were available: conditional coverage was
.9333, empirical standard deviation 9.141 and mean standard error 9.687.
Counting unavailable fits gives the previously reported .850 and .875.
The latter case therefore does not mainly show underestimated standard
errors among the usable fits. Availability, tail error and bias are different
questions. Eighty repetitions leave substantial Monte Carlo uncertainty.

The binding coordinates are transition logits, including near-zero
destinations under privately selected actions. Larger emission calibration
alone does not supply those scarce action/state transitions. The model's
positive population transition probabilities do not prevent an observed
likelihood maximum near the boundary. The observed fit reconstructions,
larger-trajectory follow-up and complete-state frequency experiment separate
this scarcity explanation from the earlier differencing/start checks. They
do not certify a global maximum or establish one universal cause of
undercoverage. At 3,000 observational trajectories and 600 independent
emission labels, all 120 follow-up fits are available; coverage is .9417,
bias $-.213$, empirical standard deviation 4.097 and mean reported standard
error 4.301. This supports a bounded larger-sample regime on this mechanism,
not a universal 3,000-customer rule. The failed 199-refit bootstrap remains
adverse evidence for the original design.

\subsection{A measurement acquisition alternative}
An independent representative audit now follows whole customer trajectories,
recording the state before each action, the same production signal, the
actual action and the eventually settled reward. The states must have an
operational definition: for example, a predecision classification of verified
repayment capacity and available liquidity into solvent, temporarily
constrained and persistent-shortfall categories. Definition and thresholds
are fixed before outcomes. Independent adjudication checks timestamped
income, obligations and accessible balances; disagreements and unavailable
documents remain measurement failures.

These categories are not automatically sufficient causal states. The model
asserts that, conditional on the state, the officer's private action
information adds no further outcome/transition cause; that the state is
Markov for this horizon; and that the audit and production emissions agree.
A coarse affordability label can violate all three. Audit selection must
represent the target population, or known inclusion probabilities and their
influence must enter the estimator. The implementation below assumes
representative independent trajectories. It does not silently replace a
convenience audit by a random sample.

Let $C_{i,sj}^E$ count signal $j$ in state $s$ along audited path $i$,
$C_{i,asj}^T$ count transitions from $s$ to $j$ under action $a$, and
$N_{i,as},R_{i,as}$ be state/action exposure and reward sums. The row-frequency
and reward estimators are
\[
 \hat E_{sj}=\frac{\sum_i C_{i,sj}^E}{\sum_{i,j'}C_{i,sj'}^E},\quad
 \hat T_{asj}=\frac{\sum_i C_{i,asj}^T}{\sum_{i,j'}C_{i,asj'}^T},\quad
 \hat r_{as}=\frac{\sum_iR_{i,as}}{\sum_iN_{i,as}}.
\]
Initial frequencies estimate $\mu$; action frequencies estimate the private
logger $q$. This is the complete-data likelihood estimate in the declared
finite model. It has no local mixture-likelihood optimization. It buys that
simplicity with a much more expensive observation design.

For any row with path counts $C_{ij}$, let $c_j=\E C_{ij}>0$ and use
logit coordinates $\ell_j=\log(c_j/c_K)$. Direct differentiation gives
\[
 \varphi_{ij}^{\ell}=C_{ij}/c_j-C_{iK}/c_K,\qquad
 \varphi_{i,as}^{r}=(R_{i,as}-r_{as}N_{i,as})/\E N_{i,as}.
\]
Stack these with the initial and emission coordinates. Their path-level
outer products preserve all within-trajectory covariance; treating visits
as independent would not. For the scaled reward coordinates divide the
reward influence by the known reward standard deviation. With fixed horizon,
positive row masses, finite second moments and accurate representative
labels, the multivariate iid central limit theorem and quotient delta method
give
\[
 \sqrt n(\hat\vartheta-\vartheta_0)
 =n^{-1/2}\sum_i\varphi_i+o_p(1),\qquad
 \widehat{\Var}(\hat V)=
 \nabla V(\hat\vartheta)'\left(n^{-2}\sum_i\hat\varphi_i\hat\varphi_i'\right)
 \nabla V(\hat\vartheta).
\]
An empty category makes these interior coordinates unavailable. Arbitrary
pseudocounts would define a different finite procedure; they are not added
to disguise support failure. This is a pointwise regular inference route,
not uniform coverage near vanishing state/action probabilities.

\subsection{A dynamic observable policy}
For a prior belief $b$ immediately before signal $o$, update
$b_o(s)=b(s)E_{so}/\sum_jb(j)E_{jo}$. If action $a$ is intervened on, the
next prior is $b_oT_a$. There is no private logger factor in that update:
the policy chose the action. For remaining horizon $h$,
\begin{equation}
 V_h(b)=\sum_oP(o\mid b)\max_a
       \{b_or_a+\gamma V_{h-1}(b_oT_a)\},\qquad V_0=0.
 \label{s6:belief-policy}
\end{equation}
The finite recursion is the belief-state planning construction
\citep{kaelbling1998}. Four signals, two actions and four decisions permit
exact tree enumeration here. The tree is optimized using an independent
2,000-trajectory audit pilot, then frozen in its entirety. Evaluation fits
kernels on separate audited trajectories and propagates their uncertainty
through this fixed tree. It differentiates neither the argmax nor the pilot
training procedure. The interval targets the value of the policy that was
actually trained, conditional on that pilot, not the unknown optimal value.

To evaluate the fixed tree without numerical posterior divisions, propagate
unnormalized state mass. At each signal branch multiply by $E_{\cdot o}$,
add the selected action's expected reward, and propagate through $T_a$.
Summing all leaves gives a polynomial in the kernels and initial law, linear
in the rewards. Its derivative is therefore well-defined for the fixed tree;
central differencing is checked against a smaller step and independent
likelihood/count perturbations. Exact enumeration avoids a filter
approximation error in this selected problem. Larger horizons require
pruning or approximate planning and an explicit value-loss bound.

On the common mechanism the frozen learned policy has true four-period
contrast 25.029 against never assisting. Some actions depend on earlier
signals even after holding the current signal and remaining horizon fixed;
this is a belief-history policy, rather than an immediate-reward rule.
At 1,200 audited evaluation trajectories, all 400 fits are available and
coverage is .950 for this contrast. The original eight-period fixed-signal
contrast has coverage .9425. At 300 audited paths, 54 of 400 fits have an
empty category; all-request coverage is .8275 for the learned policy even
though coverage conditional on availability is .9566. Measurement-supported
identification therefore still requires enough informative visits.

\subsection{Three information routes on the same servicing problem}
An observed-history g-formula substitutes the factual conditional law given
public signals and actions. Under private logging its state posterior after
action $a$ is proportional to $b_o(s)q_s(a)$, while the intervention posterior
is $b_o(s)$. Thus even an exactly estimated public-history conditional law
uses the wrong state mixture for the intervention. Adding a flexible learner
cannot remove this information gap. Recording the action-relevant private
information, or randomizing assistance with adequate support, changes that
conclusion. The common-mechanism comparison evaluates both recursions on the
same frozen policy. Its exact public-history g-formula contrast is $-35.647$,
against the intervention contrast $25.029$. The sign reversal occurs even
without nuisance estimation error.

Validated emissions can identify the kernels using the earlier marked
moment/hidden-likelihood route, with fewer labels but difficult finite
curvature and boundary behavior. Fully audited trajectories identify the
same kernels and policy target more directly, with the explicit acquisition
cost above. Their precision and availability are properties of different
sampling designs, not proof that hidden-state modeling is unnecessary.

Proximal evaluation instead seeks policy value without recovering named
states. For this servicing mechanism an action-side proxy might be a
separately recorded officer assessment, and an outcome-side proxy an
independent affordability measurement. Merely calling them proxies proves
nothing: after conditioning on the latent confounder and allowed context,
the action-side proxy and logged action must be independent of the
outcome-side proxy, proposed action and counterfactual continuation under
the source's hybrid laws. The proposed policy cannot reuse forbidden private
proxy information. \citet[Assumptions 1--2, Theorem 2 and Table 1]{bennettprox2023}
give these policy-dependent restrictions and the recursive action/outcome
bridges. In particular their outcome bridge targets
$\E[\mathbf1\{E_t=A_t\}Y_t\mid U_t,A_t=a]$, where $E_t$ denotes the
proposed action; dropping that indicator changes the theorem.

The source's recursive bridge score is a substantive challenger when the
bank can defend those exclusions and bridge existence/rank conditions but
cannot obtain accurate state labels. The production stream in this constructed servicing experiment has one
public signal stream and no independently validated second proxy. It cannot
invoke that theorem by relabelling the same signal twice. If an additional
proxy is collected, its conditional operator, range, regularization and
policy information set become new scientific inputs. This explains why the
audited trajectory route is selected for the bounded decision capability
and proximal evaluation remains a qualified alternative, rather than
claiming that one existing dataset satisfies all three information regimes.

\subsection{The initial product boundary}
The offering includes finite-state estimation and observable-history policy
evaluation under a qualified measurement design, plus the earlier exact
linear-Gaussian controlled likelihood. It does not include arbitrary nonlinear
private-state fitting with routinely calibrated policy intervals. Choosing
an extended Kalman or particle filter would require a declared likelihood,
initial/noise laws, derivative accuracy and policy-value comparison against
a numerical reference, following the approximation obligations already
derived. Section~\ref{c7:nonlinear} supplies a selected audited nonlinear
comparison with additional information and policy-class restrictions.
Those restrictions affect which bank requests can be accepted; they are
not postponed as a generic deployment check. A noisy or selected audit
also changes the statistical model, even when all matrices are well-conditioned.
In a separately registered 200-repetition sensitivity run, independently
misclassifying 10\% of predecision state labels at 1,200 audit trajectories
produces bias $-105.72$ for the original $59.56$ contrast and zero coverage.
This is a specific severe misspecification example, not a calibrated bound
for every 10\% audit-error process.


% SOURCE docs/release/2026-10-29/source/docs/writing-repair/2026-10-29/src/deb405da50a73a86.tex CHAR 0
\section{An audited nonlinear servicing experiment}
\label{c7:nonlinear}
The exact Gaussian construction leaves a substantive question: what changes
when the state transition bends and a private binary decision reveals hidden
state information? A small servicing experiment gives a complete answer for
one identified model. It also exposes the difference between an accurate
likelihood and a reliable routine interval.

For each independent customer the order is state, signal, action, reward,
then next state. In units fixed by an accurate state audit, let
\begin{align}
 S_0&\sim N(\mu_0,v_0), & O_t&=S_t+e_t,\\
 S_{t+1}&=\rho S_t+\gamma\tanh(S_t)+\beta A_t+u_{t+1},\nonumber\\
 \Pr(A_t=1\mid S_t,O_t)&=\epsilon+(1-2\epsilon)
       \operatorname{expit}(\lambda_0+\lambda_sS_t+\lambda_oO_t),\nonumber\\
 R_t&=\eta_0+\eta_s S_t+\eta_a A_t+\eta_{sa}S_tA_t+v_t^R.
 \label{c7:nonlinear-law}
\end{align}
Independent Gaussian innovations have positive variances $q,r_o,r_R$.
Their log variances, both initial parameters, three transition coefficients,
three logger coefficients and four reward coefficients give 15 unknowns.
The fixed logger floor is $\epsilon=.05$; stability requires
$|\rho|+|\gamma|<1$. Rewards are in thousands of currency units. These are
synthetic scientific restrictions, not a fitted description of a real bank.

\subsection{Identification comes from the observation design}
Production observes $(O,A,R)$; an independent representative audit observes
complete labelled $(S,O,A,R)$ trajectories. The loading one in $O=S+e$ and
the labelled scale of $S$ fix state meaning. Under positive definite population
Gram matrices, regressing $S_{t+1}$ on $(S_t,\tanh S_t,A_t)$ identifies the
three transition coefficients and its residual identifies $q$. The signal
residual identifies $r_o$. Regressing reward on $(1,S_t,A_t,S_tA_t)$ identifies
its coefficients and variance; labelled initial states identify $(\mu_0,v_0)$.
The clipped logistic probability is strictly increasing in its index. Equality
of these probabilities on a full-rank support for $(1,S,O)$ therefore implies
equality of the logger coefficients. The positive action floor supports both
actions; full rank and representative labels remain substantive requirements.
This establishes identification before filtering. Filtering does not create
the audited states. Noisy or selected labels change this argument.

Write $f_p,g_p,a_p,h_p$ for transition, signal, logger and reward densities.
A production trajectory contributes
\begin{equation}
 L_i^O(p)=\int\nu_p(s_0)
  \prod_{t=0}^{H-1}g_p(o_t\mid s_t)a_p(a_t\mid s_t,o_t)h_p(r_t\mid s_t,a_t)
  \prod_{t=0}^{H-2}f_p(s_{t+1}\mid s_t,a_t)\,ds_{0:H-1}.
 \label{c7:nonlinear-likelihood}
\end{equation}
An audited trajectory contributes this integrand evaluated at the labelled
states. Joint inference sums log likelihoods over both independent samples.
The historical logger remains in the observed likelihood; the policy replaces
it in the interventional law. This combines the marginal likelihood recursion
of \citet[Theorem 12.1]{sarkka2013} with the mechanism-replacement principle
of \citet[Corollary 1]{pearl2009} for the specified controlled model.

\subsection{Likelihood, score and value need separate numerical checks}
Take grid points $s_j$, integration weights $w_j$, initial masses
$b_{0j}=\nu_p(s_j)w_j$, emissions
$e_{tj}=g_p(o_t\mid s_j)a_p(a_t\mid s_j,o_t)h_p(r_t\mid s_j,a_t)$,
and matrices $K_{a,jk}=f_p(s_k\mid s_j,a)w_k$. Then
\[
 z_t=\sum_j b_{tj}e_{tj},\quad
 \widetilde b_{tj}=b_{tj}e_{tj}/z_t,\quad
 b_{t+1,k}=\sum_j\widetilde b_{tj}K_{a_t,jk},\quad
 \ell_i^O=\sum_t\log z_t.
\]
Finite-domain transitions are not renormalized to conceal tail loss. Scaling
emissions to prevent underflow leaves the discretized log likelihood intact.
For a parameter direction, product and quotient rules give
\begin{align*}
 \dot z_t&=\sum_j(\dot b_{tj}e_{tj}+b_{tj}\dot e_{tj}),
 &\dot\ell_t&=\dot z_t/z_t,\\
 \dot{\widetilde b}_{tj}&=(\dot b_{tj}e_{tj}+b_{tj}\dot e_{tj})/z_t
                         -\widetilde b_{tj}\dot z_t/z_t,\\
 \dot b_{t+1,k}&=\sum_j(\dot{\widetilde b}_{tj}K_{a_t,jk}
                            +\widetilde b_{tj}\dot K_{a_t,jk}).
\end{align*}
Automatic differentiation evaluates these pure array derivatives in double
precision; installed source and version identities are preserved
\citep{jaxruntime2026}. It differentiates the discretized objective, not the
unapproximated integral. Central differences in all 15 directions check the
implemented score independently.

The reference compares 41, 61, 81 and 121 nodes and domains $[-7,7]$ and
$[-9,9]$. The registered 61-node likelihood, score and value checks pass.
Every final interior fit also compares with a finer, wider grid; an interval
is withheld if the observed log-likelihood discrepancy exceeds $10^{-4}$,
the maximum score discrepancy exceeds $10^{-3}$, or value changes by more
than 0.1 currency unit. These are observed refinements, not certified error
bounds for all parameters.

The assumed-density Gaussian challenger integrates the private action by
Gauss--Hermite nodes and replaces the posterior by its first two moments.
Agreement in likelihood alone does not qualify its score or policy value.
The exact Gaussian reduction requires both $\gamma=0$ and $\lambda_s=0$:
after observing the signal, a logger depending only on that signal adds no
hidden-state information. Retaining a state-dependent logistic logger would
invalidate the reduction. The bootstrap particle challenger instead keeps
the private logger in its weights. Eight independent seeds at 512 and 4,096
particles expose Monte Carlo variation. Coupled finite differences assess
stochastic likelihood perturbations; they are not a certified derivative
through a discontinuous resampling path. These are substantive alternatives
to Gaussian projection, as emphasized by \citet[Chapters 6--7]{sarkka2013}.

\subsection{Joint audit uncertainty and the policy target}
Let $s_i^O,s_j^A$ be production and audit trajectory scores and $J$ the
negative joint Hessian. The covariance estimate is
\[
 \widehat V_p=J^{-1}\left\{
  \sum_i\widetilde s_i^O\widetilde s_i^{O\prime}+
  \sum_j\widetilde s_j^A\widetilde s_j^{A\prime}\right\}J^{-1},
\]
with within-sample centering. Sample sizes enter as sums: an audit-identified
direction does not automatically inherit the production sample's information
rate. At an interior regular solution, value variance is
$\nabla V(p)'\widehat V_p\nabla V(p)$. Boundary solutions and singular
information do not receive this interval. Forty retained full refits resample
whole trajectories separately within the two samples on the first data set.
This is a bootstrap diagnostic, not precise tail calibration. Multiple starts
and their objective differences are recorded without claiming a global optimum.

The evaluated rule is $A_t=1\{O_t>c_t\}$. An independent 800-trajectory
pilot chooses from a 64-element three-date threshold library; evaluation freezes
the selected rule. This is a permitted observable-policy subclass, not the
unknown optimal belief-history policy. At state $s$ its action probability is
$\pi_t(1\mid s)=\Phi((s-c_t)/\sqrt{r_o})$. Interventional grid mass evolves
as $b_{t+1}=(b_t\odot(1-\pi_t))K_0+(b_t\odot\pi_t)K_1$, giving
\[
 V_p(c)=\sum_{t=0}^{H-1}.97^t\sum_j b_{tj}
   \{\eta_0+\eta_s s_j+(\eta_a+\eta_{sa}s_j)\pi_t(1\mid s_j)\}.
\]
Independent policy simulation checks the value. Full enumeration benchmarks
the Gaussian planner on the same finite library. Equality of selected policies
in the retained example is a result, not general planning optimality.

On a bounded numerical domain with reward bound $R_*$, telescoping transition
errors yields a value bound of the form
$\sum_t .97^t\{\varepsilon_r+2R_*\sum_{j\le t}\varepsilon_{K,j}\}$.
The original Gaussian law is unbounded; truncation adds
$\sum_t .97^t\E[|R_t|1\{\text{path exits domain}\}]$, which cannot be
silently set to zero. Locally, mean-score error $\varepsilon_s$ and Hessian
curvature at least $\kappa$ bound parameter displacement by
$\varepsilon_s/\kappa$ within that basin. A value-gradient bound $G$ adds
$G\varepsilon_s/\kappa$ to the monetary error budget. The executed grid and
seed comparisons measure errors; they do not establish those uniform bounds.

\subsection{Boundary failures and a fresh information comparison}
The first registered baseline uses 160 audits and 80 production trajectories,
with $|\rho|\le.55$, $|\gamma|\le.4$. Only 55 of 100 fits are interior;
45 hit a boundary. Of the 55 available intervals, 54 cover. Reporting only
$54/55$ would hide availability: coverage over requested calculations is
$54/100$. Misspecified transitions yield 25 unavailable fits out of 25.
Weak signals, scarce actions, biased audits and a six-date horizon are separate
retained arms; successful optimization does not validate their assumptions.

A fresh registered comparison keeps the original policy and tests
$|\rho|+|\gamma|\le.95$, each coefficient in $[-.9,.9]$, against the original
rectangle. Other bounds stay fixed. With 160 audits the stable domain gives
99 available intervals and 94 covering draws out of 100. With 640 audits all
100 are available and 95 cover, with Monte Carlo interval $[.8872,.9836]$.
Mean standard error falls from about 168.06 to 95.84 currency units. The
rectangular challenger is available in 59 and 79 draws on these same designs.
No result switches methods after observing its interval.

This supplies an identified and executed nonlinear reference, including a
useful larger-audit regime and explicit failures. It does not establish
routine validity for arbitrary private-state laws, erroneous audits, weak
information, unrestricted policies or unknown global optimizer behavior.
Those require renewed scientific arguments and evidence.

